\section{Methods}
\label{sec:methods}

All models share the same high-level structure: an \emph{encoder} that maps each raw
sequence to a fixed-length representation, a \emph{fusion} module that combines the
enhancer and promoter representations into a joint feature vector, and a
\emph{classification head} that maps the joint vector to a scalar interaction probability.
\Cref{fig:overview} sketches this architecture.

% ---- Architecture overview figure ----
\begin{figure}[H]
\centering
\resizebox{\linewidth}{!}{%
\begin{tikzpicture}[
    node distance = 0.9cm and 1.5cm,
    box/.style    = {rectangle, draw, rounded corners, minimum width=2.2cm,
                     minimum height=0.8cm, align=center, font=\small},
    arrow/.style  = {-Latex, thick},
  ]
  \node[box, fill=blue!10]  (en_seq)  {Enhancer\\Sequence};
  \node[box, fill=blue!10,  below=of en_seq] (pr_seq) {Promoter\\Sequence};

  \node[box, fill=orange!20, right=of en_seq] (en_enc) {Encoder\\$f_\theta$};
  \node[box, fill=orange!20, right=of pr_seq] (pr_enc) {Encoder\\$f_\theta$};

  \node[box, fill=green!20,  right=2.0cm of en_enc, yshift=-0.45cm]
        (fusion) {Fusion\\Module};

  \node[box, fill=red!15,    right=of fusion]
        (head)   {Classif.\\Head};

  \node[box, fill=gray!15,   right=of head]
        (out)    {$\hat{y}$};

  \draw[arrow] (en_seq) -- (en_enc) node[midway, above, font=\scriptsize]{$s_e$};
  \draw[arrow] (pr_seq) -- (pr_enc) node[midway, above, font=\scriptsize]{$s_p$};
  \draw[arrow] (en_enc.east) -- ++(0.4,0) |- (fusion.west)
               node[near end, above, font=\scriptsize]{$\mathbf{h}_e$};
  \draw[arrow] (pr_enc.east) -- ++(0.4,0) |- (fusion.west)
               node[near end, below, font=\scriptsize]{$\mathbf{h}_p$};
  \draw[arrow] (fusion) -- (head) node[midway, above, font=\scriptsize]{$\mathbf{m}$};
  \draw[arrow] (head)   -- (out);
\end{tikzpicture}}%
\caption{High-level architecture of DeepChrInteract-v2.
         Weights of the encoder $f_\theta$ may be shared between the enhancer and
         promoter branches (parameter tying) or independent (dual-branch).}
\label{fig:overview}
\end{figure}

\subsection{Encoder Architectures}

\subsubsection{M1 / M2: CNN Baseline (Single- and Dual-Branch)}

The CNN encoder processes one-hot input $\mathbf{X} \in \mathbb{R}^{5 \times L}$
through two convolutional stages:

\begin{equation}
  \mathbf{H} = \mathrm{MaxPool} \Bigl(
    \mathrm{BN}\bigl(\mathrm{ReLU}(\mathrm{Conv1d}_{128})\bigr)
  \Bigr) \circ
  \mathrm{MaxPool} \Bigl(
    \mathrm{BN}\bigl(\mathrm{ReLU}(\mathrm{Conv1d}_{64})\bigr)
  \Bigr)(\mathbf{X}),
\end{equation}

where each stage comprises three stacked convolutions with kernel size 24 and stride 4,
followed by Batch Normalisation~\cite{batch_norm} and max-pooling.
The output is globally average-pooled and passed through two fully-connected layers
with Dropout~\cite{dropout}($p=0.5$), yielding $\mathbf{h} \in \mathbb{R}^{512}$.

\textbf{M1 (single-branch)} concatenates the enhancer and promoter sequences
along the length axis before the encoder.
\textbf{M2 (dual-branch)} runs two independent encoder copies and fuses the
resulting vectors $\mathbf{h}_e$ and $\mathbf{h}_p$ via the fusion module.
M1 is a direct reproduction of \texttt{model\_onehot\_cnn\_one\_branch} from the original codebase;
M2 reproduces \texttt{model\_onehot\_cnn\_two\_branch}.

\subsubsection{M3: \texorpdfstring{$k$}{k}-mer Embedding + CNN}

The $k$-mer tokenised sequence $\mathbf{t} \in \mathbb{Z}^{L-5}$ is first mapped through
a trainable embedding table $\mathbf{E} \in \mathbb{R}^{4097 \times 128}$:
\begin{equation}
  \hat{\mathbf{X}} = \mathbf{E}[\mathbf{t}] \in \mathbb{R}^{(L-5) \times 128}.
\end{equation}
Four stacked Conv1d--ReLU--MaxPool blocks (64 channels, kernel 32, stride 8)
progressively reduce the sequence length before a flatten and two dense layers.
This reproduces \texttt{model\_embedding\_cnn\_one/two\_branch}.

\subsubsection{M4: Bidirectional LSTM}

The BiLSTM encoder~\cite{lstm} processes the one-hot (or embedded) sequence
in both temporal directions:
\begin{equation}
  \mathbf{h}_e = \bigl[ \overrightarrow{\mathrm{LSTM}}(\mathbf{X})_{L};\;
                        \overleftarrow{\mathrm{LSTM}}(\mathbf{X})_{1} \bigr]
  \in \mathbb{R}^{2d_h},
\end{equation}
where $d_h = 256$ is the hidden size per direction and $[\cdot;\cdot]$ denotes
concatenation.  We stack two BiLSTM layers with inter-layer dropout 0.3.
The final hidden state is used as the sequence representation.

\subsubsection{M5: Transformer Encoder}

Raw one-hot sequences of length $L = 10{,}001$ are too long for standard
$O(L^2)$ attention.  We first apply a strided CNN front-end (stride 16)
to produce a downsampled sequence of $\approx 625$ tokens, then add sinusoidal
positional encoding and pass through four Transformer encoder layers
($d_{\mathrm{model}} = 256$, $n_{\mathrm{head}} = 8$,
$d_{\mathrm{ff}} = 1024$, dropout 0.1)~\cite{attention_all_you_need}.
A prepended learnable \texttt{[CLS]} token aggregates global context;
its final representation is used as the sequence vector.

\subsubsection{M6: CNN + BiLSTM Hybrid}

The CNN front-end (stride-reduced to $\approx 500$ tokens) serves as a motif
detector, and the BiLSTM captures long-range dependencies over the resulting
feature map.  This hybrid design is motivated by the finding of Singh~et~al.~\cite{speid}
that sequential models benefit from local feature extraction prior to recurrence.

\subsubsection{M7: CNN + Transformer Hybrid}

Identical to M6 but with the BiLSTM replaced by the Transformer encoder described
in M5.  The CNN front-end reduces the sequence length to a manageable scale,
after which self-attention operates without significant memory overhead.

\subsubsection{M8: DNA LLM Encoder (DNABERT / NT / HyenaDNA)}

For each of the four foundation models, we extract a fixed-length representation
$\mathbf{z} \in \mathbb{R}^{d_{\mathrm{LLM}}}$ as described in \Cref{sec:data}.
A linear projection $\mathbf{W} \in \mathbb{R}^{512 \times d_{\mathrm{LLM}}}$
aligns the LLM dimension to the common head input size.
We experiment with two modes:
\emph{frozen} (LLM parameters are fixed, only the projection and head are trained)
and \emph{fine-tuned} (all parameters are updated with a lower learning rate
of $5 \times 10^{-6}$ for the LLM layers).

\subsection{Fusion Module}
\label{sec:fusion}

Given enhancer representation $\mathbf{h}_e \in \mathbb{R}^d$ and promoter representation
$\mathbf{h}_p \in \mathbb{R}^d$, the fusion module computes an interaction vector
$\mathbf{m}$ according to one of the following strategies:

\begin{table}[H]
\centering
\caption{Supported fusion strategies. $d_{\mathrm{out}}$ is the output dimension fed
         to the classification head.}
\label{tab:fusion}
\begin{tabular}{llcc}
\toprule
Strategy & Formula & $d_{\mathrm{out}}$ & Learnable params \\
\midrule
Concat        & $[\mathbf{h}_e;\mathbf{h}_p]$               & $2d$ & 0 \\
Add           & $\mathbf{h}_e + \mathbf{h}_p$               & $d$  & 0 \\
Subtract      & $\mathbf{h}_e - \mathbf{h}_p$               & $d$  & 0 \\
Multiply      & $\mathbf{h}_e \odot \mathbf{h}_p$           & $d$  & 0 \\
Bilinear      & $\mathbf{h}_e^\top \mathbf{W}_b \mathbf{h}_p$  & 1    & $d^2$ \\
Concat+Sub+Mul & $[\mathbf{h}_e;\mathbf{h}_p;\mathbf{h}_e-\mathbf{h}_p;\mathbf{h}_e\odot\mathbf{h}_p]$ & $4d$ & 0 \\
\bottomrule
\end{tabular}
\end{table}

The \textsc{Concat+Sub+Mul} strategy is our recommended default, as it provides
a richer feature space at no additional learnable parameters, making it invariant
to the sign of the enhancer--promoter difference while preserving asymmetry
information through both the subtract and multiply paths.

\subsection{Classification Head}

All models share a two-layer MLP head:
\begin{equation}
  \hat{y} = \sigma\!\Bigl(\mathbf{W}_2\, \mathrm{ReLU}(\mathbf{W}_1 \mathbf{m} + \mathbf{b}_1) + \mathbf{b}_2\Bigr),
\end{equation}
where $\mathbf{W}_1 \in \mathbb{R}^{256 \times d_{\mathrm{out}}}$,
$\mathbf{W}_2 \in \mathbb{R}^{1 \times 256}$, and $\sigma$ is the sigmoid function.
The model is trained with binary cross-entropy loss and the Adam
optimiser~\cite{adam} ($\beta_1=0.9$, $\beta_2=0.999$, weight decay $10^{-5}$).
